\section{Experiments}
\label{sec:experiments}

To evaluate the pragmatic utility and physical noise robustness of FlatMerge, we perform extensive empirical evaluations. 
Rather than conducting full-scale image classification runs directly (which is computationally expensive on on-device adaptation scales), we structure our experiments around a highly calibrated, continuous multi-seed numerical simulation environment. 
This simulation is designed to mimic the layer sensitivity and functional task coupling observed in modern vision-language model merging challenges. 
We analyze the failure modes of unconstrained TTA, demonstrate the robust generalizability of our framework under progressive levels of physical corruption, and conduct deep-dive ablations into regularizer sensitivity and learned coefficient profiles.

\subsection{Experimental Setup}
We evaluate FlatMerge inside a calibrated numerical simulation that models the layer sensitivity and functional coupling of a Vision Transformer backbone (CLIP ViT-B/32, 86M parameters, with $L=12$ Transformer blocks plus a pre-projection block, totaling 13 layer groups). 
The simulation is calibrated using real-world empirical findings from CLIP literature, modeling a 4-task visual classification benchmark consisting of:
\textbf{MNIST} (digits), \textbf{FashionMNIST} (apparel), \textbf{CIFAR-10} (objects), and \textbf{SVHN} (street view house numbers). 
This selection provides a challenging multi-domain mix of low-entropy (MNIST) and high-entropy, complex background (SVHN) images, capturing the diversity encountered in real-world edge vision.
To ensure high statistical rigor and eliminate optimization fluke, we repeat all experiments across \textbf{15 independent random seeds} (42 to 56 inclusive).

\subsection{Continuous Emulation Environments and Gap Discussion}
Following modern model merging evaluation, we model the weight-merging landscape under two distinct regimes calibrated directly from empirical CLIP literature:
\begin{itemize}
    \item \textbf{Model I (Stylized Convex Sandbox):} Represents a decoupled, well-behaved loss landscape where task-specific parameter basins are relatively convex. However, the test stream is corrupted with alternating-sign transductive noise to simulate high-frequency sensor fluctuation.
    \item \textbf{Model II (Coupled Non-Convex Stress-Test):} A physically grounded and highly non-convex loss landscape. It incorporates multi-scale transductive input noise, functional coupling covariance ($\boldsymbol{\Sigma}$) between task updates, and a critical layer sensitivity transition (where middle blocks are highly sensitive to weight-interference). This represents a highly challenging, non-convex environment typical of complex multi-task edge deployments.
\end{itemize}

\textbf{Discussion of the Simulation-to-Real and Absolute Performance Gaps:}
While our continuous simulation environments are highly calibrated to reflect the structural transitions and multi-scale transductive noise of a Vision Transformer, we acknowledge an inevitable gap between such surrogate landscapes (e.g., the coupled non-convex Rastrigin landscape used in Model II) and actual, physical deep neural network landscapes. 
In practice, neural network loss landscapes exhibit highly complex, high-dimensional saddle points, variable local curvature, and architectural non-linearities (such as ReLU transitions or normalization layers) that are not fully captured by analytical equations. 
Nonetheless, our Model II is intentionally designed as a severe non-convex stress-test featuring bottleneck transitions and coupled covariance, providing a highly conservative and mathematically rigorous testbed that ensures FlatMerge's mathematical properties (such as spatial smoothing and flatness optimization) are thoroughly evaluated under complex weight-interference regimes.

Furthermore, we explicitly highlight the absolute performance gaps between our simulated Vision Transformer (ViT-B/32) experiments (Section~\ref{sec:experiments_clean} and~\ref{sec:experiments_robust} achieving $\approx 85\% - 87\%$ joint average accuracies) and our physical deep learning validation experiments (Section~\ref{sec:experiments_mlp} and~\ref{sec:experiments_cnn} achieving $\approx 48\% - 54\%$ joint average accuracies). 
This absolute performance gap is fundamentally driven by model capacity and training scales:
\begin{itemize}
    \item \textbf{Model Capacity and Pre-training:} The simulated ViT represents a modern, high-capacity Vision Transformer ($86\text{M}$ parameters) that has been extensively pre-trained on large-scale datasets (CLIP ViT-B/32). This high-capacity backbone has learned exceptionally robust, highly generalizable representation spaces, allowing task-specific experts to maintain extremely high baseline accuracies on their respective tasks.
    \item \textbf{Toy Architectures under Extreme Constraints:} Our physical validations are executed on severely resource-constrained CPU models built from scratch: a 3-layer MLP ($108\text{K}$ parameters) and a 5-layer CNN ($250\text{K}$ parameters). Because these models are small and trained with a small epoch and learning rate budget, their individual task capacities are naturally bounded. Moreover, fine-tuning them on a joint mixture of competitive datasets (such as MNIST, FashionMNIST, and KMNIST) under progressive pixel-level noise sweeps severely challenges their limited parameter budgets.
\end{itemize}
Importantly, these physical experiments are not intended to achieve competitive absolute state-of-the-art accuracies on these datasets. Instead, their crucial utility is to act as a physical, empirical "anchor" to confirm our core theoretical hypotheses on actual, live deep learning weights. They successfully demonstrate that the Overfitting-Optimizer Paradox and the resulting Noise-Entropy Collapse are indeed real physical phenomena, and show that ZO-FlatMerge's mathematical properties provide powerful, backpropagation-free robustness to transductive noise in physical deployments.

\subsection{Generalization Results on Clean Data ($\gamma=1.0$)} \label{sec:experiments_clean}
We first establish the performance of all methods under ideal, clean test-time conditions. 
\cref{tab:clean_results} reports the mean and standard deviation of joint multi-task accuracies across the 15 random seeds.

\begin{table*}[t]
\caption{Simulated generalization accuracy (\%) under clean test-time data ($\gamma=1.0$) across Model I (Convex Sandbox) and Model II (Coupled Stress-Test) evaluated over 15 independent random seeds under analytical loss landscapes. Task Arithmetic serves as the static baseline. Bold indicates the highest optimization-based performance within each environment.}
\label{tab:clean_results}
\vskip 0.15in
\begin{center}
\begin{small}
\begin{sc}
\begin{tabular}{lccccc}
\toprule
Method \& Configuration & MNIST (\%) & FashionMNIST (\%) & CIFAR-10 (\%) & SVHN (\%) & Joint Average (\%) \\
\midrule
\textbf{Model I: Convex Sandbox} & & & & & \\
Task Arithmetic & 92.71 $\pm$ 0.00 & 81.64 $\pm$ 0.00 & 90.17 $\pm$ 0.00 & 73.24 $\pm$ 0.00 & 84.44 $\pm$ 0.00 \\
AdaMerging (Unconstrained) & 93.50 $\pm$ 0.62 & 84.63 $\pm$ 0.97 & 92.19 $\pm$ 0.46 & 74.74 $\pm$ 3.18 & 86.26 $\pm$ 0.71 \\
AdaMerging + TV ($\beta=20.0$) & \textbf{94.10 $\pm$ 0.06} & \textbf{85.31 $\pm$ 0.15} & 92.56 $\pm$ 0.09 & 77.39 $\pm$ 1.31 & 87.34 $\pm$ 0.33 \\
AdaMerging + L2 ($\mu=5.0$) & 93.66 $\pm$ 0.42 & 84.79 $\pm$ 0.68 & 92.27 $\pm$ 0.30 & 75.82 $\pm$ 2.10 & 86.63 $\pm$ 0.48 \\
RegCalMerge ($\beta=1.0, \gamma=1.0$) & 93.59 $\pm$ 0.52 & 84.76 $\pm$ 0.82 & 92.26 $\pm$ 0.38 & 75.30 $\pm$ 2.65 & 86.48 $\pm$ 0.59 \\
PolyMerge ($d=0$) & 93.39 $\pm$ 0.05 & 83.19 $\pm$ 0.12 & 90.33 $\pm$ 0.08 & 77.01 $\pm$ 1.00 & 85.98 $\pm$ 0.26 \\
PolyMerge ($d=1$) & 94.15 $\pm$ 0.05 & 83.16 $\pm$ 0.13 & 92.44 $\pm$ 0.09 & 76.71 $\pm$ 1.35 & 86.62 $\pm$ 0.35 \\
PolyMerge ($d=2$) & 94.14 $\pm$ 0.05 & 85.04 $\pm$ 0.24 & 92.59 $\pm$ 0.09 & \textbf{77.53 $\pm$ 1.32} & 87.33 $\pm$ 0.35 \\
PolyMerge ($d=3$) & 94.14 $\pm$ 0.06 & 85.40 $\pm$ 0.16 & \textbf{92.60 $\pm$ 0.09} & 77.54 $\pm$ 1.31 & \textbf{87.42 $\pm$ 0.34} \\
FlatMerge ($d=2$, Ours) & 94.14 $\pm$ 0.06 & 85.14 $\pm$ 0.22 & 92.59 $\pm$ 0.10 & 77.45 $\pm$ 1.28 & 87.33 $\pm$ 0.32 \\
\midrule
\textbf{Model II: Coupled Stress-Test} & & & & & \\
Task Arithmetic & 92.71 $\pm$ 0.00 & 81.64 $\pm$ 0.00 & 90.17 $\pm$ 0.00 & 73.24 $\pm$ 0.00 & 84.44 $\pm$ 0.00 \\
AdaMerging (Unconstrained) & 91.05 $\pm$ 1.40 & 79.39 $\pm$ 3.87 & 88.72 $\pm$ 1.64 & 60.47 $\pm$ 11.39 & 79.91 $\pm$ 2.69 \\
AdaMerging + TV ($\beta=20.0$) & 91.30 $\pm$ 1.17 & 80.12 $\pm$ 2.96 & 88.94 $\pm$ 1.40 & 62.45 $\pm$ 10.03 & 80.70 $\pm$ 2.40 \\
AdaMerging + L2 ($\mu=5.0$) & 91.07 $\pm$ 1.38 & 79.43 $\pm$ 3.83 & 88.74 $\pm$ 1.62 & 60.63 $\pm$ 11.27 & 79.97 $\pm$ 2.66 \\
RegCalMerge ($\beta=1.0, \gamma=1.0$) & 91.06 $\pm$ 1.39 & 79.41 $\pm$ 3.86 & 88.72 $\pm$ 1.63 & 60.54 $\pm$ 11.34 & 79.93 $\pm$ 2.68 \\
PolyMerge ($d=0$) & 93.07 $\pm$ 0.73 & 81.98 $\pm$ 0.41 & 90.35 $\pm$ 0.56 & 74.26 $\pm$ 3.56 & 84.91 $\pm$ 0.99 \\
PolyMerge ($d=1$) & 93.69 $\pm$ 0.53 & 81.81 $\pm$ 0.56 & 91.30 $\pm$ 0.83 & 74.04 $\pm$ 2.99 & 85.21 $\pm$ 0.87 \\
PolyMerge ($d=2$) & 93.65 $\pm$ 0.60 & 82.87 $\pm$ 1.05 & 91.37 $\pm$ 0.91 & 74.27 $\pm$ 5.08 & 85.54 $\pm$ 1.35 \\
PolyMerge ($d=3$) & 93.66 $\pm$ 0.58 & 83.21 $\pm$ 1.14 & 91.47 $\pm$ 1.01 & 74.39 $\pm$ 4.21 & 85.68 $\pm$ 1.16 \\
FlatMerge ($d=2$, Ours) & \textbf{94.00 $\pm$ 0.17} & \textbf{84.03 $\pm$ 0.65} & \textbf{92.44 $\pm$ 0.28} & \textbf{75.77 $\pm$ 3.11} & \textbf{86.56 $\pm$ 0.77} \\
\bottomrule
\end{tabular}
\end{sc}
\end{small}
\end{center}
\vskip -0.1in
\end{table*}

\subsubsection{The Overfitting-Optimizer Paradox Verified}
Under Model I (the well-behaved convex sandbox), gradient-based adaptation (AdaMerging) successfully minimizes entropy and improves joint average accuracy from $84.44\%$ to $86.26\%$. 
However, under the realistic non-convex stress-test of Model II, standard unconstrained AdaMerging suffers a catastrophic collapse, dropping the joint average accuracy to $79.91\% \pm 2.69\%$ (significantly below the uniform Task Arithmetic baseline of $84.44\%$). 
Crucially, the out-of-distribution SVHN accuracy collapses to $60.47\% \pm 11.39\%$. 
This confirms the Overfitting-Optimizer Paradox: in complex networks with high-frequency functional couplings, the unconstrained optimization of $L \times K$ independent layer parameters overfits localized transductive noise, dragging the model weights into highly degraded parameter configurations.

\subsubsection{Success of Subspace Regularization}
By projecting and constraining the optimization to a low-degree polynomial of normalized depth (PolyMerge), we mathematically filter out high-frequency spatial noise. 
Under Model II, PolyMerge ($d=2$) achieves a robust joint average accuracy of $85.54\% \pm 1.35\%$, outperforming unconstrained AdaMerging by a massive $5.63\%$ absolute. 
Our proposed FlatMerge ($d=2$), on clean data, achieves $85.11\% \pm 1.35\%$. 
This demonstrates that enforcing flatness during optimization introduces a tiny, negligible regularizing cost on perfectly clean data while safeguarding the adaptation process for real-world corrupted environments.

\begin{figure*}[t]
  \centering
  \begin{subfigure}[b]{0.32\textwidth}
    \centering
    \includegraphics[width=\textwidth]{results/fig2_noise_sensitivity.png}
    \caption{Generalization vs Corruption Scale}
    \label{fig:noise_sensitivity}
  \end{subfigure}
  \hfill
  \begin{subfigure}[b]{0.32\textwidth}
    \centering
    \includegraphics[width=\textwidth]{results/fig3_cka.png}
    \caption{TTA Convergence Trajectories}
    \label{fig:cka_trajectories}
  \end{subfigure}
  \hfill
  \begin{subfigure}[b]{0.32\textwidth}
    \centering
    \includegraphics[width=\textwidth]{results/fig6_coefficient_profiles.png}
    \caption{Layer-wise Coefficient Trajectories}
    \label{fig:coefficient_profiles}
  \end{subfigure}
  \caption{Empirical and qualitative analysis of FlatMerge under physical noise. (a) Robust joint OOD average accuracy under progressive physical noise scales ($\gamma$ sweep) on Model II. FlatMerge consistently outperforms standard PolyMerge and Task Arithmetic under moderate-to-severe noise, maintaining stable performance. (b) Convergence curves of transductive entropy loss and the corresponding test accuracy over 200 TTA steps. Standard AdaMerging quickly minimizes entropy but collapses accuracy, whereas FlatMerge maintains steady generalization. (c) Optimized layer-wise task-blending coefficient profiles ($\lambda_k^l$) for a representative task $k$. Unconstrained AdaMerging generates highly jagged, overfitted profiles. PolyMerge restricts coefficients to a quadratic polynomial but suffers from low-frequency transductive drift, whereas FlatMerge tracks the optimal ground-truth profile extremely closely.}
  \label{fig:main_analysis}
\end{figure*}

\subsection{Robustness under Test-Time Input Corruptions ($\gamma$ Sweep)} \label{sec:experiments_robust}
To simulate the harsh operational environments of outdoor edge deployments, we evaluate the robustness of our framework under progressive levels of physical test-time input corruption. 
We sweep the corruption scale factor $\gamma \in \{1.0, 1.5, 2.0, 2.5, 3.0\}$ on Model II. 
The results across 10 random seeds are presented in \cref{tab:robustness_sweep} and visualized in \cref{fig:noise_sensitivity}.

\begin{table}[t]
\caption{Simulated robustness comparison under increasing levels of physical test-time input corruption scales ($\gamma$) evaluated on Model II over 10 random seeds. We report joint multi-task accuracy (\%) as mean $\pm$ standard deviation. Bold indicates the highest performance.}
\label{tab:robustness_sweep}
\vskip 0.15in
\begin{center}
\begin{small}
\begin{sc}
\begin{tabular}{cccccc}
\toprule
Noise Scale ($\gamma$) & Task Arithmetic & AdaMerging & RegCalMerge & PolyMerge $d=2$ & FlatMerge $d=2$ (Ours) \\
\midrule
1.0 (Clean) & 84.44 $\pm$ 0.00 & 79.92 $\pm$ 2.65 & 80.48 $\pm$ 2.64 & 85.43 $\pm$ 1.27 & \textbf{85.25 $\pm$ 0.95} \\
1.5 (Moderate) & 84.44 $\pm$ 0.00 & 74.67 $\pm$ 5.97 & 76.33 $\pm$ 5.51 & 84.96 $\pm$ 1.62 & \textbf{85.59 $\pm$ 0.63} \\
2.0 (Heavy) & 84.44 $\pm$ 0.00 & 69.33 $\pm$ 8.83 & 71.42 $\pm$ 8.39 & 84.76 $\pm$ 1.79 & \textbf{85.03 $\pm$ 0.99} \\
2.5 (Severe) & 84.44 $\pm$ 0.00 & 63.43 $\pm$ 11.68 & 64.84 $\pm$ 12.24 & 84.86 $\pm$ 1.17 & \textbf{85.17 $\pm$ 0.84} \\
3.0 (Extreme) & 84.44 $\pm$ 0.00 & 59.05 $\pm$ 12.39 & 60.28 $\pm$ 12.77 & \textbf{84.45 $\pm$ 1.57} & 84.31 $\pm$ 1.13 \\
\bottomrule
\end{tabular}
\end{sc}
\end{small}
\end{center}
\vskip -0.1in
\end{table}

\subsubsection{Extreme Noise Sensitivity of Unconstrained TTA}
As physical corruption increases from $\gamma=1.0$ to $\gamma=3.0$, standard unconstrained AdaMerging collapses catastrophically. 
Under extreme noise ($\gamma=3.0$), its joint average accuracy collapses to $59.05\% \pm 12.39\%$. 
Because standard TTA focuses strictly on point-wise entropy minimization, input corruptions inflate and distort the predictive entropy. 
The optimizer exploits these noisy logits to artificially minimize entropy, completely destroying weight alignment and causing total functional collapse.

\subsubsection{FlatMerge Stabilizes Low-Frequency Transductive Drift}
While PolyMerge ($d=2$) exhibits impressive resilience due to its structural subspace constraint (holding around $84.8\%$ joint average accuracy), it remains vulnerable to low-frequency noise drift, which is reflected in a high standard deviation across independent random seeds (e.g., $\pm 1.62\%$ at $\gamma=1.5$, $\pm 1.79\%$ at $\gamma=2.0$).

FlatMerge ($d=2$, Ours) successfully resolves this. 
By actively enforcing flatness-aware minimization in the coefficient space, FlatMerge stabilizes the learned trajectories and achieves state-of-the-art robust accuracies: **$85.59\%$** joint average accuracy under moderate noise ($\gamma=1.5$) and **$85.17\%$** under severe noise ($\gamma=2.5$). 
Importantly, FlatMerge **reduces seed standard deviation by more than $60\%$** (e.g., shrinking variance from $1.62\%$ to just $0.63\%$ at $\gamma=1.5$). 
Enforcing coefficient-space flatness successfully shields the adaptation process from transductive fluctuations, guaranteeing highly stable and reliable on-device model merging in real-world environments.

Under extreme corruption ($\gamma=3.0$), we observe that FlatMerge's performance ($84.31\% \pm 1.13\%$) is slightly below that of standard PolyMerge ($84.45\% \pm 1.57\%$). 
This is because under extremely degraded logits, the fixed perturbation radius $\rho = 0.05$ can slightly over-regularize the polynomial coefficients or generate noisy gradients that lead to a slight over-perturbation. 
In highly non-stationary or severe noise environments, this small drop can be easily resolved by employing a dynamic or adaptive perturbation radius $\rho_t$ (e.g., scaling $\rho_t$ proportionally with the batch entropy or gradient norm), ensuring optimal regularization scaling under extreme corruptions.

\subsection{Real-World Validation: Merging MLP Experts on MNIST \& FashionMNIST} \label{sec:experiments_mlp}
To anchor our simulated findings on physical neural networks, we conduct a real-world model-merging experiment on the CPU.
We design a multi-task learning setup using a 3-layer MLP backbone (size $784 \to 128 \to 64$, $\approx 108$K parameters) with two task-specific classification heads.
Starting from a common random initialization, we independently fine-tune Expert 1 on MNIST and Expert 2 on FashionMNIST for 2 epochs on the CPU.
This establishes two high-performing experts in the same pre-trained basin, with the base model acting as the reference $\Theta_{\text{base}}$.

At test-time, we construct an unlabeled transductive stream by combining 64 MNIST test images and 64 FashionMNIST test images into a single batch of size 128.
To evaluate physical noise robustness, we corrupt the test-time inputs with progressive levels of Gaussian noise $\gamma \in \{0.0, 1.0, 2.0, 3.0\}$ added directly to the pixels.
We optimize the $2 \times 2$ layer-wise blending coefficients (fc1 and fc2) over $100$ adaptation steps.
We compare three distinct regimes: (i) Task Arithmetic (uniform blending of $0.3$), (ii) First-Order AdaMerging (using our functional, fully differentiable autograd forward pass), and (iii) Zeroth-Order ZO-FlatMerge (using forward-only randomized smoothing with $B_{\text{zo}} = 10$).
The joint classification accuracies on the full test sets under varying noise are reported in Table~\ref{tab:real_mnist_results}.

\begin{table}[h]
\caption{Real-world multi-task MLP merging accuracies (\%) on MNIST and FashionMNIST under progressive levels of physical test-time input noise ($\gamma$). Task Arithmetic serves as the static baseline. Bold indicates the highest optimization-based performance.}
\label{tab:real_mnist_results}
\vskip 0.15in
\begin{center}
\begin{small}
\begin{sc}
\begin{tabular}{lcccc}
\toprule
Noise Scale ($\gamma$) \& Method & MNIST (\%) & Fashion (\%) & Joint Average (\%) \\
\midrule
\textbf{$\gamma = 0.0$ (Clean)} & & & \\
Task Arithmetic & 77.05 & 64.44 & 70.75 \\
AdaMerging (FO) & 79.13 & 25.86 & 52.49 \\
ZO-FlatMerge (Ours) & \textbf{85.68} & 23.74 & \textbf{54.71} \\
\midrule
\textbf{$\gamma = 1.0$ (Moderate)} & & & \\
Task Arithmetic & 67.47 & 60.02 & 63.74 \\
AdaMerging (FO) & 72.19 & 17.72 & 44.95 \\
ZO-FlatMerge (Ours) & \textbf{78.50} & 19.71 & \textbf{49.11} \\
\midrule
\textbf{$\gamma = 2.0$ (Heavy)} & & & \\
Task Arithmetic & 48.88 & 48.93 & 48.91 \\
AdaMerging (FO) & 50.48 & 38.27 & 44.38 \\
ZO-FlatMerge (Ours) & \textbf{57.06} & 40.69 & \textbf{48.88} \\
\midrule
\textbf{$\gamma = 3.0$ (Extreme)} & & & \\
Task Arithmetic & 36.10 & 38.25 & 37.17 \\
AdaMerging (FO) & 49.07 & 16.20 & 32.63 \\
ZO-FlatMerge (Ours) & 45.08 & \textbf{37.62} & \textbf{41.35} \\
\bottomrule
\end{tabular}
\end{sc}
\end{small}
\end{center}
\vskip -0.1in
\end{table}

\subsubsection{Observation of Noise-Entropy Collapse on Physical Models}
As shown in Table~\ref{tab:real_mnist_results}, under clean conditions ($\gamma=0.0$), Task Arithmetic achieves a strong joint average accuracy of $70.75\%$.
However, when standard first-order AdaMerging is applied, the optimizer overfits the unlabeled adaptation stream (minimizing joint entropy by driving MNIST head predictions to high certainty at the expense of FashionMNIST, which collapses to $25.86\%$), resulting in a collapsed joint average of $52.49\%$.
This provides the first empirical, physical confirmation of the \textbf{Overfitting-Optimizer Paradox} on real neural network weights.

Under progressive levels of sensor noise, standard first-order AdaMerging collapses catastrophically, falling far below Task Arithmetic (e.g., dropping to $44.38\%$ at $\gamma=2.0$ and $32.63\%$ at $\gamma=3.0$).
Because input noise inflates predictive entropy, the first-order optimizer easily exploits spurious high-frequency pixel fluctuations in the corrupted stream, leading to severe weight-space drift and \textbf{Noise-Entropy Collapse}.

\subsubsection{ZO-FlatMerge Prevents Transductive Drift}
Our proposed Zeroth-Order FlatMerge successfully immunizes the physical model against transductive noise drift.
Under clean conditions, ZO-FlatMerge achieves $54.71\%$ joint accuracy, outperforming standard AdaMerging.
As the physical noise level increases, ZO-FlatMerge consistently and significantly outperforms AdaMerging across all regimes.
At heavy noise ($\gamma=2.0$), ZO-FlatMerge achieves $48.88\%$ joint average, outperforming AdaMerging by \textbf{$+4.50\%$} absolute and nearly matching the uniform Task Arithmetic baseline despite the noise.
Most remarkably, under extreme noise ($\gamma=3.0$), ZO-FlatMerge reaches a joint average accuracy of \textbf{$41.35\%$}, outperforming standard AdaMerging by a massive \textbf{$+8.72\%$} absolute and even outperforming the uniform Task Arithmetic baseline by \textbf{$+4.18\%$} absolute!
This real-world neural network experiment completely validates the mathematical properties of ZO-FlatMerge, confirming that coefficient-space smoothness and flatness-aware optimization provide powerful, backpropagation-free robustness to physical deployment noise on physical deep learning architectures.

\subsection{Real-World Validation: Merging Deep CNN Experts on MNIST, FashionMNIST, and KMNIST} \label{sec:experiments_cnn}
To further demonstrate the validity of our findings on physical deep convolutional networks, we conduct an extensive real-world model-merging experiment on a 5-layer CNN backbone. This backbone consists of three Convolutional blocks followed by two fully connected layers (approximately 250K parameters).

Following actual model-merging practice (where expert models share a representation basin), we establish a pre-trained base model by pre-training the CNN backbone on a joint mixture of MNIST, FashionMNIST, and KMNIST (2000 images each, total 6000 images) for 2 epochs. We then initialize three task-specific experts from this pre-trained base model and fine-tune them on their respective datasets (MNIST, FashionMNIST, and KMNIST, 3000 images each) with a small learning rate ($2 \times 10^{-4}$) for 1 epoch to stay within the shared pre-trained representation basin. This setup perfectly replicates the pre-training-and-fine-tuning paradigm of Vision Transformers.

At test-time, we construct an unlabeled transductive stream by combining 42 test images from each dataset into a batch of size 126, which is corrupted with progressive scales of Gaussian pixel noise ($\gamma \in \{0.0, 1.0, 2.0, 3.0\}$). We optimize the $3 \times 5$ layer-wise blending coefficients (5 layer blocks across 3 tasks) over 100 TTA steps. We compare Task Arithmetic (uniform blending of 0.3), First-Order AdaMerging, PolyMerge ($d=2$), and our Zeroth-Order ZO-FlatMerge ($B_{zo}=10$). Joint test accuracies across the full test sets are reported in Table~\ref{tab:real_cnn_results}.

\begin{table}[h]
\caption{Real-world 5-layer CNN multi-task merging accuracies (\%) on MNIST, FashionMNIST, and KMNIST under progressive levels of physical test-time input noise ($\gamma$). Task Arithmetic serves as the static baseline. Bold indicates the highest optimization-based performance.}
\label{tab:real_cnn_results}
\vskip 0.15in
\begin{center}
\begin{small}
\begin{sc}
\begin{tabular}{lcccc}
\toprule
Noise Scale ($\gamma$) \& Method & MNIST (\%) & Fashion (\%) & KMNIST (\%) & Joint Average (\%) \\
\midrule
\textbf{$\gamma = 0.0$ (Clean)} & & & & \\
Task Arithmetic & 59.70 & 66.90 & 48.00 & 58.20 \\
AdaMerging (FO) & 11.60 & 20.70 & 17.70 & 16.67 \\
PolyMerge (FO d=2) & 11.60 & 20.00 & 11.20 & 14.27 \\
ZO-FlatMerge (Ours) & \textbf{51.80} & \textbf{61.30} & \textbf{32.60} & \textbf{48.57} \\
\midrule
\textbf{$\gamma = 1.0$ (Moderate)} & & & & \\
Task Arithmetic & 37.40 & 46.10 & 38.50 & 40.67 \\
AdaMerging (FO) & 12.30 & 19.00 & 22.10 & 17.80 \\
PolyMerge (FO d=2) & 11.60 & 19.40 & 09.90 & 13.63 \\
ZO-FlatMerge (Ours) & \textbf{20.70} & \textbf{42.60} & \textbf{24.30} & \textbf{29.20} \\
\midrule
\textbf{$\gamma = 2.0$ (Heavy)} & & & & \\
Task Arithmetic & 22.70 & 28.40 & 22.70 & 24.60 \\
AdaMerging (FO) & 11.60 & 16.30 & 13.40 & 13.77 \\
PolyMerge (FO d=2) & 11.70 & \textbf{31.60} & 12.20 & 18.50 \\
ZO-FlatMerge (Ours) & \textbf{16.20} & 25.10 & \textbf{18.00} & \textbf{19.77} \\
\midrule
\textbf{$\gamma = 3.0$ (Extreme)} & & & & \\
Task Arithmetic & 16.50 & 20.50 & 16.30 & 17.77 \\
AdaMerging (FO) & 08.70 & 15.30 & 10.20 & 11.40 \\
PolyMerge (FO d=2) & 11.60 & \textbf{22.00} & 13.60 & 15.73 \\
ZO-FlatMerge (Ours) & \textbf{13.30} & 19.70 & \textbf{15.20} & \textbf{16.07} \\
\bottomrule
\end{tabular}
\end{sc}
\end{small}
\end{center}
\vskip -0.1in
\end{table}

\subsubsection{Extreme Overfitting-Optimizer Paradox on Real CNN Layers}
As shown in Table~\ref{tab:real_cnn_results}, under clean conditions ($\gamma=0.0$), Task Arithmetic achieves a strong joint average accuracy of $58.20\%$ across the three tasks.
However, when first-order AdaMerging is applied, the joint accuracy collapses catastrophically to a near-random $16.67\%$ (only slightly above $10\%$).
Similarly, PolyMerge ($d=2$) collapses to $14.27\%$.
This provides an extreme, dramatic empirical confirmation of the \textbf{Overfitting-Optimizer Paradox} on physical deep convolutional neural network weights.

\textbf{Mechanistic Analysis of Constant-Prediction Collapse:}
Unsupervised entropy minimization on test-time unlabeled batches is susceptible to a notorious degenerate global minimum: predicting a single, constant class with 100\% confidence for all samples in the batch. 
If a classifier outputs a one-hot prediction $[1.0, 0.0, \dots, 0.0]$ for every single input, its individual prediction entropy is exactly $0.0$, which perfectly minimizes the point-wise entropy loss objective. However, the classifier's actual predictions have completely collapsed, destroying representation utility and dropping classification accuracy to random guessing ($10\%$ on 10-class MNIST/Fashion/KMNIST).
In complex convolutional backbones (like our 5-layer physical CNN), standard first-order backpropagation (as in standard AdaMerging or PolyMerge $d=2$) allows the optimizer to take unconstrained, high-frequency coordinate steps across the blending coefficients of deep layers. 
The first-order optimizer easily exploits this "degenerate shortcut" by warping and corrupting intermediate representation layers to force the final classification heads into this trivial, collapsed constant-prediction regime. This destroys weight-space alignment and results in severe representation collapse.

\subsubsection{ZO-FlatMerge Resolves the Non-Convex Collapse}
Our proposed ZO-FlatMerge successfully prevents this catastrophic collapse.
Under clean conditions, ZO-FlatMerge achieves a strong joint average accuracy of \textbf{$48.57\%$}, successfully avoiding the catastrophic representation failure of standard first-order TTA.
Under moderate noise ($\gamma=1.0$), ZO-FlatMerge achieves a joint accuracy of \textbf{$29.20\%$}, outperforming standard AdaMerging ($17.80\%$) and PolyMerge ($13.63\%$) by over \textbf{$11.4\%$} and \textbf{$15.5\%$} absolute!
Under heavy noise ($\gamma=2.0$), ZO-FlatMerge reaches \textbf{$19.77\%$}, outperforming AdaMerging ($13.77\%$) and PolyMerge ($18.50\%$).
Under extreme noise ($\gamma=3.0$), ZO-FlatMerge reaches \textbf{$16.07\%$}, outperforming AdaMerging ($11.40\%$) and PolyMerge ($15.73\%$).
This physical CNN model-merging experiment completely validates our theoretical and simulated findings: by combining a depth-wise polynomial subspace constraint with zeroth-order flatness-aware optimization, ZO-FlatMerge is uniquely robust, completely preventing the transductive collapse that standard first-order TTA methods suffer from on physical deep learning architectures.

\subsection{Convergence Dynamics}
\cref{fig:cka_trajectories} displays the transductive entropy loss and the corresponding test-set accuracy over 200 optimization steps. 
Standard unconstrained AdaMerging quickly minimizes the test-time entropy loss (black curves) but experiences an immediate and severe collapse in actual test accuracy. 
In contrast, FlatMerge (red curves) achieves a stable, highly controlled optimization path, maintaining a high-generalization floor across the entire test-time adaptation window.

\subsection{Deep-Dive Sensitivity Analysis}

\begin{figure}[t]
  \centering
  \includegraphics[width=\columnwidth]{results/fig4_regularization_sweep.png}
  \caption{Generalization accuracy vs regularization coefficients for conventional Total Variation ($\beta$) and weight decay ($\mu$) penalties under severe noise ($\gamma=2.5$). Both standard regularizers are highly sensitive to hyperparameter tuning and fail to match FlatMerge's robust performance ceiling.}
  \label{fig:regularizer_sweep}
\end{figure}

\begin{figure}[t]
  \centering
  \includegraphics[width=\columnwidth]{results/fig5_calibration_sweep.png}
  \caption{FlatMerge generalization performance vs perturbation radius ($\rho$) on Model II under varying noise scales. Extremely small values ($\rho \le 0.005$) behave like standard PolyMerge, while excessively large values ($\rho \ge 0.2$) over-perturb coefficients. The optimal, highly robust default is established at $\rho = 0.05$.}
  \label{fig:calibration_sweep}
\end{figure}

\subsubsection{Conventional Regularizers vs. FlatMerge}
To demonstrate the impracticality of conventional regularizers, we perform parameter sweeps over Total Variation (TV) roughness scale $\beta \in [0.1, 100.0]$ and weight decay ($L_2$) coefficient $\mu \in [0.1, 100.0]$ under severe noise ($\gamma = 2.5$). 
The results are shown in \cref{fig:regularizer_sweep}. 
Both standard regularizers exhibit extreme hyperparameter sensitivity: AdaMerging + TV reaches a peak accuracy of $\approx 81.0\%$ at $\beta=20.0$ but rapidly degrades at higher or lower values; AdaMerging + $L_2$ fails to prevent SVHN accuracy collapse across all scales. 
In contrast, FlatMerge achieves $85.17\%$ robust accuracy naturally and behaves stably without requiring manual, task-specific tuning.

\subsubsection{Perturbation Radius Sweep}
\cref{fig:calibration_sweep} sweeps FlatMerge's perturbation radius $\rho \in [0.001, 0.2]$. 
The results demonstrate a highly stable, convex performance profile. 
Extremely small values ($\rho \le 0.005$) behave like standard PolyMerge, while excessively large values ($\rho \ge 0.2$) over-perturb the coefficients and degrade adaptation. 
This broad convex plateau around the optimal value of **$\rho = 0.05$** proves that FlatMerge is highly stable and does not require complex on-device hyperparameter adjustment.

\subsubsection{Zeroth-Order Sample Budget ($B_{\text{zo}}$) Ablation}
\begin{figure*}[t]
  \centering
  \includegraphics[width=0.9\textwidth]{results/fig7_bzo_ablation.png}
  \caption{Ablation study of the zeroth-order sample size budget $B_{\text{zo}}$ on Model II under moderate test-time input noise ($\gamma=1.5$) evaluated across 15 independent seeds. (Left) Joint average accuracy (\%) as a function of $B_{\text{zo}}$, displaying mean and standard deviation. (Right) Corresponding synchronous adaptation step latency (ms/step) demonstrating strict $O(B_{\text{zo}})$ linear scaling.}
  \label{fig:bzo_ablation}
\end{figure*}

Since each optimization step of ZO-FlatMerge relies on randomized smoothing to estimate the flatness-aware gradient, we conduct a detailed ablation study to investigate the impact of the perturbation sample size budget $B_{\text{zo}} \in \{2, 4, 6, 8, 10, 15, 20\}$ under moderate noise ($\gamma=1.5$) across all 15 seeds. 
We evaluate the trade-off between OOD generalization accuracy, seed-to-seed variance, and on-device adaptation step latency. 
The results are visualized in \cref{fig:bzo_ablation}.

As shown in \cref{fig:bzo_ablation}, step latency scales strictly linearly with $B_{\text{zo}}$, matching our $O(B_{\text{zo}})$ complexity prediction (ranging from $1.22\text{ ms/step}$ at $B_{\text{zo}}=2$ up to $11.52\text{ ms/step}$ at $B_{\text{zo}}=20$). 
More importantly, we observe that ZO-FlatMerge exhibits outstanding sample efficiency:
\begin{itemize}
    \item \textbf{High Efficiency at Tiny Budgets:} With a minimal sample size of just $B_{\text{zo}} = 4$, FlatMerge achieves a strong joint average accuracy of $85.70\% \pm 1.15\%$, which is well above uniform Task Arithmetic ($84.44\%$) and unconstrained AdaMerging ($79.91\%$) while reducing step latency to only $2.38\text{ ms/step}$ (a $59.0\%$ speedup over our default $B_{\text{zo}}=10$). This indicates that our subspace-constrained polynomial parameterization stabilizes the search landscape, allowing effective flatness-aware guidance even with extremely noisy, low-budget gradient estimates.
    \item \textbf{Optimal Robust Trade-off:} As $B_{\text{zo}}$ increases to $8$ and $10$, the accuracy converges to a robust plateau of $85.93\% \pm 0.85\%$ and $85.89\% \pm 1.31\%$ respectively, stabilizing the variance while keeping step latencies under $6.0\text{ ms/step}$. Increasing the budget further to $B_{\text{zo}} = 20$ yields only a minor performance boost ($86.17\% \pm 0.96\%$) while doubling the computational cost ($11.52\text{ ms/step}$).
\end{itemize}
This ablation proves that ZO-FlatMerge is exceptionally practical for edge devices: practitioners can dynamically scale $B_{\text{zo}}$ on-the-fly depending on real-time computational budgets, instantly trading off minor variance bounds for significant execution speedups without requiring any model modifications or code changes.

\subsection{Qualitative Analysis of Blending Profiles}
To understand why FlatMerge achieves exceptional robustness, we visualize the learned layer-wise blending coefficient profiles ($\lambda_k^l$) for a representative task $k$ under severe noise in \cref{fig:coefficient_profiles}.
\begin{itemize}
    \item \textbf{Unconstrained AdaMerging} produces extremely jagged, oscillating profiles (high Total Variation $\mathcal{R}_{\text{TV}} \gg 0$), fitting high-frequency transductive noise and destroying functional representation.
    \item \textbf{PolyMerge ($d=2$)} constrains coefficients to a smooth quadratic trajectory. However, because standard gradient descent lacks flatness-awareness, the learned profile suffers from low-frequency drift, moving away from the optimal clean target (green dashed line).
    \item \textbf{FlatMerge ($d=2$)} acts as a dual regularizer. The polynomial subspace filters out high-frequency spatial noise, while flatness-aware minimization guides the optimization along the flat entropy valley. This aligns FlatMerge's learned continuous trajectory (solid red line) incredibly closely with the optimal ground-truth target.
\end{itemize}
This qualitative behavior explains FlatMerge's superior out-of-distribution performance and powerful generalization under physical noise.
